
\chapter{BNPL: кредит как платёжный метод}\label{ch:bnpl}
BNPL (Buy Now Pay Later) на уровне протокола -- точка входа
в~кредитный продукт. Со стороны клиента это выглядит как рассрочка:
выбрал товар, разбил оплату на части, забрал покупку сразу. В~момент
оплаты BNPL-провайдер принимает кредитное решение в~реальном времени:
оценивает клиента, одобряет или отклоняет заявку, и~в~случае
одобрения перечисляет продавцу полную сумму покупки сразу. Клиент
затем возвращает деньги провайдеру равными частями по графику.

Продавец получает выручку так же быстро, как при обычном эквайринге,
но платит за это комиссию BNPL-провайдеру. Кредитный риск целиком
лежит на провайдере: если клиент не выплатит оставшиеся части,
продавец денег не потеряет. Распределение рисков и определяет
экономику BNPL: провайдер зарабатывает на комиссии продавца
и~иногда на пенях за просрочку, а~продавец получает рост конверсии
и увеличение среднего чека за счёт более низкого порога входа
для покупателя.

На уровне протокола существуют два паттерна.

\section{BNPL в~роли PSP}\label{sec:bnpl-as-psp}

В~этой модели BNPL-провайдер работает как отдельный PSP. ТСП
заключает с~ним \textit{merchant agreement}, встраивает его SDK или
перенаправление в~чекаут и обрабатывает обратные вызовы с~результатом
кредитного решения.

Клиент на странице оплаты выбирает
BNPL-провайдера. Провайдер показывает клиенту условия: количество
платежей, даты, сумму каждого платежа. Клиент подтверждает. Провайдер
проводит скоринг и~выносит решение в~реальном времени. При одобрении
провайдер перечисляет продавцу полную сумму покупки за вычетом комиссии.
Продавец видит одну транзакцию на полную сумму и отгружает товар.
Далее провайдер самостоятельно списывает с~клиента оставшиеся
платежи по графику -- продавец в~этом процессе не участвует.

\begin{figure}[tbp]
  \centering
  \includefiguresvg[width=\linewidth]{assets/figures/ch21-bnpl/bnpl-psp-flow}
  \caption{BNPL в~роли PSP: платёжная последовательность и распределение ролей.}\label{fig:bnpl-psp-flow}
\end{figure}

Интеграция похожа на подключение обычного PSP: API
для создания заказа, серверные уведомления (webhook) об изменении
статуса, API для возвратов. Спорные сценарии определяются договором
с~провайдером, а~не правилами платёжной системы. Карточного
возвратного платежа здесь нет: карточная сеть в~платеже не участвует.

По этой модели работают Klarna, Affirm, Tabby и~все основные
российские BNPL-сервисы.

\section{BNPL-на-карте}\label{sec:bnpl-on-card}

Второй паттерн встраивает рассрочку в~карточную инфраструктуру. Здесь
клиент расплачивается картой, а разбиение на части происходит
на стороне эмитента или отдельного кредитора внутри карточной сети.
Продавец может вообще не знать, что покупка оплачена в~рассрочку.

Visa и Mastercard реализуют это по-разному. Visa Installments
строится на наборе правил MIT (Merchant-Initiated Transaction).
Первая авторизация проходит как обычная карточная операция
с~хранимыми данными (\textit{stored credential}) и POS Environment
кодом~I (installment MIT). После неё продавец или
планировщик Visa последовательно выставляет транши по сохранённым
реквизитам. Кредитный риск лежит на эмитенте: он решает,
для каких карт и~каких покупок доступна рассрочка. ТСП должно
поддерживать Installments API у~своего процессора, а~эмитент должен
активировать программу для своих карт -- без этого предложение
рассрочки клиенту не появится~\cite{visa-installments}.

Mastercard Installments в~основной реализации работает иначе.
Участвующий кредитор -- банк или BNPL-финтех, контрактованный
с Mastercard -- выпускает одноразовый виртуальный номер карты на
отдельном BIN. Клиент расплачивается этим номером у~продавца.
Продавец получает полную сумму сразу как обычную карточную
транзакцию. Кредитор самостоятельно собирает взносы
с~клиента~\cite{mc-installments}. Продавцу не нужно
подключать отдельный API или активировать программу, однако
по выделенному BIN транзакцию можно идентифицировать.

Различие между двумя паттернами меняет интеграционные требования.
BNPL в~роли PSP требует отдельного \textit{merchant agreement} с~провайдером:
добавить его как метод оплаты без нового договора нельзя.
BNPL-на-карте требует поддержки Installments API у~процессора
(Visa) или наличия участвующего кредитора в~сети (Mastercard),
но не требует от ТСП отдельной интеграции с BNPL-провайдером.

\section{Российский BNPL}\label{sec:bnpl-russia}

Российский рынок BNPL сформировался по модели «BNPL в~роли PSP».
Отличие от западных аналогов: все крупные российские
BNPL-сервисы принадлежат банкам или финансовым группам, у~которых
уже есть банковская лицензия и кредитная инфраструктура. Сервисы
принимают карты любых банков, а~не только собственных клиентов;
ограничивает не эмитент карты, а~скоринг.

\enquote{Долями} (Т-Банк) -- первый российский BNPL-сервис,
запущенный в~апреле 2021 года~\cite{tbank-dolyame}. Продавец
интегрирует его как внешний PSP: отдельный договор, API или
плагин для CMS, отдельный поток
уведомлений~\cite{tbank-dolyame-business}. \enquote{Подели}
устроен аналогично; формальный оператор сервиса -- ООО
\enquote{А-Четыре Технологии}, 100\%-дочернее общество
АО \enquote{Альфа-Банк}; интеграция через API, плагины для CMS
или через платёжных
партнёров~\cite{alfa-podeli,alfa-podeli-business}.
\enquote{Плати частями} (Сбер) -- отдельный BNPL-сервис,
оператор которого -- ООО \enquote{Центр новых финансовых
сервисов} (ЦНФС), 100\%-дочернее общество Сбербанка. Продавцу
необходимо подключить его как самостоятельный платёжный способ
у~процессора -- он не связан с~интеграцией
SberPay~\cite{sber-pay-parts}.

Перечисленные сервисы объединяет одна архитектура: провайдер
работает как PSP с~собственным merchant agreement и скорингом,
а~продавец получает полную сумму покупки за вычетом комиссии.
Различие -- в~модели подключения. \enquote{Долями}
и \enquote{Подели} требуют отдельного договора с~оператором.
\enquote{Сплит} (Яндекс) работает внутри Yandex Pay SDK:
продавец подключает его через единую интеграцию с Yandex Pay,
не заключая отдельного BNPL-договора~\cite{yandex-split}.
\enquote{Плати частями} подключается как отдельный метод
у~процессора.

\section{Регулирование}\label{sec:bnpl-regulation}

До недавнего времени BNPL в~большинстве юрисдикций существовал
в~регуляторном пробеле: формально не кредит (процентов для клиента
нет), но по экономической сути -- кредитный продукт.

В~России этот пробел закрыт Федеральным законом от 31~июля
2025 года 283-ФЗ~\cite{fz-283}: основной массив положений
вступает в~силу 1~апреля 2026~года, отдельные нормы (в~части
требований к~ОСР по статье~11) -- с 1~января 2027~года;
конкретный график применения сверяется по тексту закона
перед публикацией. Закон вводит понятие \enquote{оператор
сервиса рассрочки} (ОСР) и устанавливает для него отдельный
регуляторный режим: регистрация в~реестре Банка
России~\cite{cbr-osr}, требования к~капиталу, обязательное
раскрытие полной стоимости рассрочки для клиента, ограничение
максимального срока и~запрет на взимание вознаграждения
с~клиента. При совокупной задолженности клиента свыше
порога, установленного 283-ФЗ (по действующей редакции
порядка 50~тыс.\ рублей; точное значение сверяется по тексту
закона), данные передаются в~бюро кредитных историй.
Провайдер, который хочет продолжать работу после вступления
в~силу основного массива закона, должен получить статус ОСР.

В~ЕС BNPL включён в~обновлённую Consumer Credit
Directive~\cite{eu-ccd-2023}: обязательная оценка
кредитоспособности, единый стандарт раскрытия условий, право
клиента на отказ.

Провайдер, не получивший статус ОСР или не прошедший
европейскую авторизацию, перестанет работать. При выборе
BNPL-провайдера проверьте его регуляторный статус на дату
интеграции.

\practicenote{%
  При выборе BNPL-провайдера проверьте долю одобрений среди вашей
  аудитории и регуляторный статус провайдера
  (для России -- наличие в~реестре ОСР).

}
